% Diagram of the proof of Theorem C down to the cost of one lift (head of Part 4).
\begin{figure}[tbp]
\centering
\begin{tikzpicture}[x=1cm,y=1cm]
% row 0
\node[pfmain] (C) at (-0.82,0) {Thm.~\ref{main:thm:C}\\Theorem~C};
% row 1
\node[pfnode, text width=1.8cm] (aux) at (-6.145,-1.8) {Thm.~\ref{det-chain:thm:aux}\\auxiliary forms};
\node[pfnode, text width=1.8cm] (height) at (-0.82,-1.8) {Thm.~\ref{det-chain:thm:height}\\the inverse ratio};
\node[pfnode, text width=3.0cm] (twl) at (3.9,-1.8) {Thm.~\ref{setup:thm:twist}, Prop.~\ref{setup:prop:targetlattice}\\a twist for each prime; target lattice (Part~1)};
% row 2
\node[pfnode, text width=2.4cm] (chain) at (-3.64,-3.6) {Thms.~\ref{det-chain:thm:inj}--\ref{det-chain:thm:factor}\\injectivity, norm, factorisation};
\node[pfpart, text width=1.8cm] (regular) at (-0.82,-3.6) {Thm.~\ref{cg-selmer:thm:regular}\\regular families\\Fig.~\ref{fig:part4b}};
\node[pfnode, text width=2.3cm] (olmain) at (2.3,-3.6) {Thm.~\ref{ol-values:thm:main}\\the cost of one lift};
% row 3
\node[pfnode, text width=2.1cm] (onelift) at (-4.945,-5.6) {Prop.~\ref{det-chain:prop:onelift}\\norm of one lift};
\node[pfnode, text width=2.1cm] (triangle) at (-2.335,-5.6) {Lem.~\ref{det-chain:lem:triangle}, Thm.~\ref{det-chain:thm:triangle}\\triangle of the chain};
\node[pfnode, text width=2.1cm] (comp) at (0.6,-5.6) {Thm.~\ref{ol-values:thm:comp}\\comparison operator};
\node[pfnode, text width=2.1cm] (inv) at (4.04,-5.6) {Prop.~\ref{ol-values:prop:rallis}--Cor.~\ref{ol-values:cor:hv}\\Hecke values, inverse ratios};
% row 4
\node[pfnode, text width=1.8cm] (mult) at (-6.145,-7.7) {Cor.~\ref{sp-moment:cor:mult}\\equal multiplicities (Part~3)};
\node[pffact, text width=1.8cm] (rallisdc) at (-3.835,-7.7) {Fact~\ref{det-chain:fact:rallis}\\Rallis' inner product formula};
\node[pfnode, text width=2.2cm] (olc) at (-0.68,-7.7) {\S\S\ref{ol-compare:sec:exchanges}--\ref{ol-compare:sec:span}\\three exchanges, seesaw for one lift, products span};
\node[pfnode, text width=1.8cm] (periods) at (1.83,-7.7) {Thm.~\ref{ol-values:thm:periods}\\periods of false lifts vanish};
\node[pfnode, text width=1.6cm] (sharp) at (4.04,-7.7) {Thm.~\ref{der:thm:sharp}\\sharpness (Part~2)};
\node[pffact, text width=1.8cm] (rallisol) at (6.25,-7.7) {Fact~\ref{ol-values:fact:rallis}\\Rallis' inner product formula};
% row 5
\node[pfnode, text width=1.8cm] (square) at (-2.335,-9.8) {Thm.~\ref{det-exchange:thm:square}\\one square};
\node[pfnode, text width=1.8cm] (explicit) at (0.1,-9.8) {Thm.~\ref{det-exchange:thm:explicit}\\explicit real data};
\node[pffact, text width=2.0cm] (tower) at (2.6,-9.8) {Fact~\ref{ol-values:fact:tower}\\descents, periods, Hecke operators};
% row 6
\node[pfnode, text width=2.6cm] (exchange) at (-1.2,-11.6) {Thm.~\ref{det-exchange:thm:exchange}, Props.~\ref{det-exchange:prop:top}--\ref{det-exchange:prop:equalrank}\\exchange identity};
% row 7
\node[pffact, text width=2.6cm] (factsdx) at (-1.2,-13.2) {Facts~\ref{det-exchange:fact:tower}--\ref{det-exchange:fact:lib}\\tower, constancy, three inputs};
% arrows: premise -> conclusion
\draw[pfedge] (aux) -- (C);
\draw[pfedge] (chain) -- (C);
\draw[pfedge] (height) -- (C);
\draw[pfedge] (twl) -- (C);
\draw[pfedge] (chain) -- (height);
\draw[pfedge] (regular) -- (height);
\draw[pfedge] (olmain) -- (height);
\draw[pfedge] (onelift) -- (chain);
\draw[pfedge] (triangle) -- (chain);
\draw[pfedge] (comp) -- (olmain);
\draw[pfedge] (inv) -- (olmain);
\draw[pfedge] (mult) -- (onelift);
\draw[pfedge] (rallisdc) -- (onelift);
\draw[pfedge] (rallisdc) -- (triangle);
\draw[pfedge] (square) -- (triangle);
\draw[pfedge] (olc) -- (comp);
\draw[pfedge] (periods) -- (comp);
\draw[pfedge] (periods) -- (inv);
\draw[pfedge] (sharp) -- (inv);
\draw[pfedge] (rallisol) -- (inv);
\draw[pfedge] (explicit) -- (olc);
\draw[pfedge] (tower) -- (periods);
\draw[pfedge] (exchange) -- (square);
\draw[pfedge] (exchange) -- (explicit);
\draw[pfedge] (exchange.west) -- (-7.45,-11.6) |- (aux.west);
\draw[pfedge] (factsdx) -- (exchange);
\end{tikzpicture}
\caption{The proof of Theorem~C down to the cost of one lift (Part~4,
Sections~\ref{sec:det-chain}--\ref{sec:ol-compare}). Each box is a result, with its number and a
short name; a box with several numbers, or with the number of a subsection, stands for a group of
results. An arrow from one box to
another means that the proof of the second, as written in this paper, uses the first. Dashed grey
boxes are background results admitted without proof; the double box stands for the existence of
regular families, whose proof is drawn in Figure~\ref{fig:part4b}. Boxes marked Part~1, 2 or~3
are results of the earlier parts.}
\label{fig:part4a}
\end{figure}
